% GENERATED FILE. Edit canonical lessons, then run npm run generate:books.
% canonical-source-hash: fnv1a64:96c3eb9a47466269
% canonical-lessons: BN-C510-phinail, BN-C510-jhul, BN-C510-saph, BN-C510-ningrano, BN-C510-maja

\chapter{Cleaning Day}
\label{ch:bn-cleaning-day}
\hlchaptermodality{510}

\begin{chapteropening}
\textbf{By the end of this chapter:} \emph{I can talk about cobwebs, scrubbing and wringing out cloths.}

Complete the last lesson of chapter 510: I can talk about cobwebs, scrubbing and wringing out cloths.
\end{chapteropening}

\section[\texorpdfstring{\bn{ফিনাইল} (phināil)}{ফিনাইল (phināil)}]{\bn{ফিনাইল} (phināil) — phenyl (a floor disinfectant)}

\label{lesson:BN-C510-phinail}



\begin{quote}
Before the new one: say the Bengali for a cushion, then the Bengali for to spread out.
\end{quote}

\subsection*{You'll want to know: \bn{ফিনাইল}}
\textbf{\bn{ফিনাইল}} — \emph{phināil} — \textquotedblleft{}phenyl (a floor disinfectant)\textquotedblright{}.

\begin{quote}
\bn{আমি} \bn{জলে} \bn{ফিনাইল} \bn{মিশিয়ে} \bn{মেঝে} \bn{মুছলাম।}
\end{quote}

\begin{quote}
\emph{Āmi jôle phināil mishiye mejhe muchhlām.}
\end{quote}

\begin{quote}
I mixed phenyl into the water and mopped the floor.
\end{quote}

\begin{cousinweb}
\textbf{\bn{ফিনাইল}} (\emph{phināil}) is borrowed from English \emph{phenyl}, the name of a disinfectant used for washing floors.
\end{cousinweb}

\begin{grammarlens}[title={What you've built}]
One more word for describing what you lived and felt.
\end{grammarlens}

\subsection*{Your turn}
\begin{itemize}
\raggedright
  \item \emph{Say it:} \emph{phināil}
  \item \emph{Say it:} \emph{phināil}, once more
  \item \emph{Say it:} \emph{Āmi jôle phināil mishiye mejhe muchhlām.} — I mixed phenyl into the water and mopped the floor.
  \item \emph{Recall:} say the Bengali for a cushion, then the Bengali for to spread out, then say \emph{phināil} again
\end{itemize}

\begin{tcolorbox}[breakable,colback=teal!4,colframe=teal!35!black,title={Before you move on}]
What does \bn{ফিনাইল} mean? (\textquotedblleft{}Phenyl (a floor disinfectant)\textquotedblright{}.) Say it once more.
\end{tcolorbox}
\section[\texorpdfstring{\bn{ঝুল} (jhul)}{ঝুল (jhul)}]{\bn{ঝুল} (jhul) — cobwebs; sooty dust}

\label{lesson:BN-C510-jhul}



\begin{quote}
Before the new one: say the Bengali for to spread out, then the Bengali for floor disinfectant.
\end{quote}

\subsection*{You'll want to know: \bn{ঝুল}}
\textbf{\bn{ঝুল}} — \emph{jhul} — \textquotedblleft{}cobwebs; sooty dust\textquotedblright{}.

\begin{quote}
\bn{সিলিং} \bn{থেকে} \bn{আমি} \bn{ঝুল} \bn{ঝাড়লাম।}
\end{quote}

\begin{quote}
\emph{Siling theke āmi jhul jhāṛlām.}
\end{quote}

\begin{quote}
I brushed the cobwebs off the ceiling.
\end{quote}

\begin{cousinweb}
\textbf{\bn{ঝুল}} (\emph{jhul}) probably belongs with \emph{jhulyati}, \textquotedblleft{}swings, hangs\textquotedblright{} (CDIAL 5406), the root of \bn{ঝোলানো}: the dust that hangs from the ceiling.
\end{cousinweb}

\begin{grammarlens}[title={What you've built}]
One more word for describing what you lived and felt.
\end{grammarlens}

\subsection*{Your turn}
\begin{itemize}
\raggedright
  \item \emph{Say it:} \emph{jhul}
  \item \emph{Say it:} \emph{jhul}, once more
  \item \emph{Say it:} \emph{Siling theke āmi jhul jhāṛlām.} — I brushed the cobwebs off the ceiling.
  \item \emph{Recall:} say the Bengali for to spread out, then the Bengali for floor disinfectant, then say \emph{jhul} again
\end{itemize}

\begin{tcolorbox}[breakable,colback=teal!4,colframe=teal!35!black,title={Before you move on}]
What does \bn{ঝুল} mean? (\textquotedblleft{}Cobwebs; sooty dust\textquotedblright{}.) Say it once more.
\end{tcolorbox}
\section[\texorpdfstring{\bn{সাফ} (sāph)}{সাফ (sāph)}]{\bn{সাফ} (sāph) — clean; cleared away}

\label{lesson:BN-C510-saph}



\begin{quote}
Before the new one: say the Bengali for floor disinfectant, then the Bengali for cobwebs.
\end{quote}

\subsection*{You'll want to know: \bn{সাফ}}
\textbf{\bn{সাফ}} — \emph{sāph} — \textquotedblleft{}clean; cleared away\textquotedblright{}.

\begin{quote}
\bn{ঝড়ের} \bn{পরে} \bn{আমরা} \bn{উঠোনটা} \bn{সাফ} \bn{করলাম।}
\end{quote}

\begin{quote}
\emph{Jhôṛer pôre āmrā uṭhonṭā sāph kôrlām.}
\end{quote}

\begin{quote}
After the storm we cleared the courtyard.
\end{quote}

\begin{cousinweb}
\textbf{\bn{সাফ}} (\emph{sāph}) is usually traced to Arabic \emph{ṣāf}, \textquotedblleft{}pure, clear\textquotedblright{}, through Persian.
\end{cousinweb}

\begin{grammarlens}[title={What you've built}]
One more word for describing what you lived and felt.
\end{grammarlens}

\subsection*{Your turn}
\begin{itemize}
\raggedright
  \item \emph{Say it:} \emph{sāph}
  \item \emph{Say it:} \emph{sāph}, once more
  \item \emph{Say it:} \emph{Jhôṛer pôre āmrā uṭhonṭā sāph kôrlām.} — After the storm we cleared the courtyard.
  \item \emph{Recall:} say the Bengali for floor disinfectant, then the Bengali for cobwebs, then say \emph{sāph} again
\end{itemize}

\begin{tcolorbox}[breakable,colback=teal!4,colframe=teal!35!black,title={Before you move on}]
What does \bn{সাফ} mean? (\textquotedblleft{}Clean; cleared away\textquotedblright{}.) Say it once more.
\end{tcolorbox}
\section[\texorpdfstring{\bn{নিংড়ানো} (ningṛāno)}{নিংড়ানো (ningṛāno)}]{\bn{নিংড়ানো} (ningṛāno) — to wring out}

\label{lesson:BN-C510-ningrano}



\begin{quote}
Before the new one: say the Bengali for cobwebs, then the Bengali for cleared away.
\end{quote}

\subsection*{You'll want to know: \bn{নিংড়ানো}}
\textbf{\bn{নিংড়ানো}} — \emph{ningṛāno} — \textquotedblleft{}to wring out\textquotedblright{}.

\begin{quote}
\bn{আমি} \bn{ভিজে} \bn{গামছাটা} \bn{নিংড়ালাম।}
\end{quote}

\begin{quote}
\emph{Āmi bhije gāmchhāṭā ningṛālām.}
\end{quote}

\begin{quote}
I wrung out the wet towel.
\end{quote}

Past (he/she …): \bn{তিনি} \bn{কাপড়টা} \bn{নিংড়ালেন}, \emph{tini kāpṛṭā ningṛālen}.

\begin{cousinweb}
\textbf{\bn{নিংড়ানো}} (\emph{ningṛāno}) comes from a reconstructed \emph{nirgaḍati}, \textquotedblleft{}drips out\textquotedblright{}; CDIAL lists Bengali \emph{niṅṛāna}, \textquotedblleft{}to make drip, to squeeze\textquotedblright{} (CDIAL 7293).
\end{cousinweb}

\begin{grammarlens}[title={What you've built}]
One more word for describing what you lived and felt.
\end{grammarlens}

\subsection*{Your turn}
\begin{itemize}
\raggedright
  \item \emph{Say it:} \emph{ningṛāno}
  \item \emph{Say it:} \emph{ningṛāno}, once more
  \item \emph{Say it:} \emph{Āmi bhije gāmchhāṭā ningṛālām.} — I wrung out the wet towel.
  \item \emph{Recall:} say the Bengali for cobwebs, then the Bengali for cleared away, then say \emph{ningṛāno} again
\end{itemize}

\begin{tcolorbox}[breakable,colback=teal!4,colframe=teal!35!black,title={Before you move on}]
What does \bn{নিংড়ানো} mean? (\textquotedblleft{}To wring out\textquotedblright{}.) Say it once more.
\end{tcolorbox}
\section[\texorpdfstring{\bn{মাজা} (mājā)}{মাজা (mājā)}]{\bn{মাজা} (mājā) — to scour (pots); to brush (teeth)}

\label{lesson:BN-C510-maja}



\begin{quote}
Before the new one: say the Bengali for cleared away, then the Bengali for to wring out.
\end{quote}

\subsection*{You'll want to know: \bn{মাজা}}
\textbf{\bn{মাজা}} — \emph{mājā} — \textquotedblleft{}to scour (pots); to brush (teeth)\textquotedblright{}.

\begin{quote}
\bn{রাতে} \bn{শোয়ার} \bn{আগে} \bn{আমি} \bn{দাঁত} \bn{মাজি।}
\end{quote}

\begin{quote}
\emph{Rāte shoyār āge āmi dā̃t māji.}
\end{quote}

\begin{quote}
I brush my teeth at night before bed.
\end{quote}

Past (he/she …): \bn{তিনি} \bn{কড়াইটা} \bn{মাজলেন}, \emph{tini kôṛāiṭā mājlen}.

\begin{cousinweb}
\textbf{\bn{মাজা}} (\emph{mājā}) is a tadbhava: CDIAL lists Bengali \emph{mājā}, \textquotedblleft{}to scour\textquotedblright{}, under Sanskrit \emph{mārjati}, \textquotedblleft{}wipes clean\textquotedblright{} (CDIAL 10080).
\end{cousinweb}

\begin{grammarlens}[title={What you've built}]
One more word for describing what you lived and felt.
\end{grammarlens}

\subsection*{Your turn}
\begin{itemize}
\raggedright
  \item \emph{Say it:} \emph{mājā}
  \item \emph{Say it:} \emph{mājā}, once more
  \item \emph{Say it:} \emph{Rāte shoyār āge āmi dā̃t māji.} — I brush my teeth at night before bed.
  \item \emph{Recall:} say the Bengali for cleared away, then the Bengali for to wring out, then say \emph{mājā} again
\end{itemize}

\begin{tcolorbox}[breakable,colback=teal!4,colframe=teal!35!black,title={Before you move on}]
What does \bn{মাজা} mean? (\textquotedblleft{}To scour (pots); to brush (teeth)\textquotedblright{}.) Say it once more.
\end{tcolorbox}
